%================================================================
\section{The village model and the targets}\label{sec:model}

\subsection{The village}

The world is a graph and a matrix of covariates, drawn once with fixed seeds and
shared by every simulation. The graph follows an edge rule rather than a target
mean degree: for every ordered pair the code draws
$A_{ij}\sim\mathrm{Bernoulli}(p)$ with $p=\EdgeRate/N$, symmetrizes the result
and sets the diagonal to zero. The expected mean degree is \MeanDegreeTheory{},
and the drawn graph has mean degree \MeanDegree{}. It is close to tree-like
(Appendix~\ref{app:model}). Each person has three
covariates $X_i\sim U[0,1]^3$, which stand for education, resources and a
baseline propensity.

Each person has a level of preventive behavior $Y_{it}\in(0,1)$. Starting from
$Y_i(0)\sim U[0,1]$, the kernel updates every person at once for
$t=0,\ldots,T-1$ with $T=\Horizon$:
\begin{equation}\label{eq:kernel}
  Y_i(t+1)
  \;=\;
  \Lambda\!\Bigl(
    \alpha + \beta\, Y_i(t)
    + \gamma\, \frac{(A\,Y(t))_i}{k_i}
    + \eta^\top X_i
    + \kappa\, d_i
    + \sigma_\varepsilon\,\varepsilon_i(t)
  \Bigr),
\end{equation}
where $\Lambda(u)=1/(1+e^{-u})$, $\varepsilon_i(t)\sim N(0,1)$ and seeding acts
from the first period on. The index is linear in means, passed through a
logistic link. The update rule is chosen, not derived; its lineage is the
linear-in-means model of \citet{manski1993}, the row-normalized peer mean on a
network of \citet{bramoulle2009}, and the logistic link of discrete choice with
social interactions in \citet{brock2001}. Peers enter through their lagged
outcomes, not through their seeding status, so the effect of a seed spreads along
the graph over time. This is the source of the interference that the analysts must
handle.

Five parameters are free, $c=(\alpha,\beta,\gamma,\kappa,\sigma_\varepsilon)$,
and they range over a box $C$ fixed before any run (Table~\ref{tab:box}). The
covariate weights $\eta$ are fixed inside the kernel. The box excludes $\gamma=0$
on purpose, since the project is about interference. Over the whole box the update
map is a contraction in the sup norm, so for a fixed assignment the process has a
unique stationary regime and forgets its starting state at a geometric rate
(Appendix~\ref{app:model}). The kernel has no tipping points and no multiple
equilibria; in the language of Section~\ref{sec:pi}, it is complete.

\begin{table}[t]
\centering
\caption{The free parameters, their box $C$, and the planted point $\cstar$.}
\label{tab:box}
\begin{zebratabular}{llll}
\toprule
\colhead{Parameter} & \colhead{Box} & \colhead{Planted value} & \colhead{Role} \\
\midrule
$\alpha$ & $[\BoxAlphaLo,\,\BoxAlphaHi]$ & $\CStarAlpha$ & intercept of the index \\
$\beta$ & $[\BoxBetaLo,\,\BoxBetaHi]$ & $\CStarBeta$ & weight on the own lag \\
$\gamma$ & $[\BoxGammaLo,\,\BoxGammaHi]$ & $\CStarGamma$ & weight on the peer mean \\
$\kappa$ & $[\BoxKappaLo,\,\BoxKappaHi]$ & $\CStarKappa$ & effect of being seeded \\
$\sigma_\varepsilon$ & $[\BoxSigmaLo,\,\BoxSigmaHi]$ & $\CStarSigma$ & SD of the shock \\
$\eta$ & fixed & $(\EtaOne,\,\EtaTwo,\,\EtaThree)$ & covariate weights \\
\bottomrule
\end{zebratabular}
\tablenote{The planted point is the center of $C$. The box, $\eta$, the graph
rule, $N$, $K$ and $T$ are exogenous and live in \texttt{config/world.toml}.}
\end{table}

\subsection{The targets: the status quo and one randomized pilot}\label{sec:targets}

The targets are what a researcher could observe. She sees the village as it is,
under $\pi_0$, and she can run one randomized pilot, under $\pi_1$. She never sees
the policy $\pi_2$. That policy is what she wants to predict, so it is never a
target; if it were, the targets would contain the answer to the planner's
question.\footnote{An earlier design of this project did exactly that: its targets
used outcomes under both seeding policies, and their difference was the targeting
premium. It is archived at the git tag \texttt{v1-worldbank-moments} and plays no
role here.}

The code uses $\NMoments$ moments in a fixed order (Table~\ref{tab:moments}).
Each is a mean over $\RepsMoment$ paired replications, and rep $r$ of the
status-quo arm and of the pilot arm share one state stream. The Monte Carlo
variance of each moment is the sample variance over the replications divided by
the number of finite values; the code drops non-finite correlations, and if fewer
than two values are finite the variance is infinite and the point is rejected.

\begin{table}[t]
\centering
\caption{The target moments, in the order the code uses.}
\label{tab:moments}
\begin{zebratabularx}{\linewidth}{clllX}
\toprule
\colhead{\#} & \colhead{Name} & \colhead{Arm} & \colhead{Target $z$} & \colhead{Definition} \\
\midrule
1 & \texttt{sq\_level} & status quo & $\ZOne$ & mean over people of $Y_T$ \\
2 & \texttt{sq\_corr\_x1} & status quo & $\ZTwo$ & correlation of $Y_T$ with the first covariate \\
3 & \texttt{sq\_peer\_corr} & status quo & $\ZThree$ & correlation of $Y_{iT}$ with the peer mean $(AY_T)_i/k_i$ \\
4 & \texttt{pilot\_level} & pilot & $\ZFour$ & mean over people of $Y_T$ \\
\bottomrule
\end{zebratabularx}
\tablenote{Targets are planted: $z=\widehat M(\cstar)$ with $\RepsZ$ replications,
stored in \texttt{data/planted.json}.}
\end{table}

The targets are planted: $z=\widehat M(\cstar)$, the simulated moments at the
planted point, computed with $\RepsZ$ replications together with their Monte Carlo
variances. Planting gives a known benchmark. It does not guarantee that a finite
search keeps any point, and it makes the true effect known, which is the only way
to check an interval against the truth.

Because the status-quo and pilot arms share their state streams, the difference
$\texttt{pilot\_level}-\texttt{sq\_level}$ at $\cstar$ equals $\theta_1(\cstar)$
exactly: the two printed values are $\PilotMinusSq$ and $\ThetaOneStar$. This is
intended. A randomized pilot identifies the effect of its own design. The
question of the paper is what it says about $\theta_2$, a contrast that no target
contains.
